% Gemeinsamer Kopf für Lernplan und Klausurskript.
% Gemeinsame Präambel für alle LaTeX-Dateien.
\documentclass[a4paper,11pt]{article}
\usepackage[utf8]{inputenc}
\usepackage[T1]{fontenc}
\usepackage[ngerman]{babel}
\usepackage{amsmath,amssymb,amsthm}
\usepackage{enumitem}
\usepackage{booktabs}
\setcounter{MaxMatrixCols}{30}

\theoremstyle{definition}
\newtheorem{definition}{Definition}[section]
\newtheorem{satz}[definition]{Satz}
\newtheorem{lemma}[definition]{Lemma}
\newtheorem{korollar}[definition]{Korollar}
\newtheorem{beispiel}[definition]{Beispiel}
\newtheorem{bemerkung}[definition]{Bemerkung}
\newtheorem{algorithmus}[definition]{Algorithmus}
\newtheorem{aufgabe}{Aufgabe}[section]
\newenvironment{loesung}{\par\noindent\textbf{Lösung.}\ }{\par}

% Kurzbefehle
\newcommand{\N}{\mathbb{N}}
\newcommand{\Z}{\mathbb{Z}}
\newcommand{\Q}{\mathbb{Q}}
\newcommand{\R}{\mathbb{R}}
\newcommand{\Zm}[1]{\mathbb{Z}_{#1}}
\newcommand{\ggT}{\operatorname{ggT}}
\newcommand{\kgV}{\operatorname{kgV}}

\usepackage[a4paper,margin=2cm]{geometry}
\usepackage{fancyhdr}
\usepackage{xcolor}
\usepackage[most]{tcolorbox}
\usepackage{array}
\usepackage{longtable}
\usepackage{amssymb}
\setlength{\parindent}{0pt}\setlength{\parskip}{4pt}
\newtcolorbox{ergebnis}{colback=green!6,colframe=green!45!black,boxrule=0.6pt,left=6pt,right=6pt,top=3pt,bottom=3pt,sharp corners}
\newtcolorbox{achtung}{colback=orange!8,colframe=orange!70!black,boxrule=0.6pt,left=6pt,right=6pt,top=3pt,bottom=3pt,sharp corners}
\newtcolorbox{rezept}[1]{colback=blue!4,colframe=blue!50!black,boxrule=0.6pt,left=6pt,right=6pt,top=3pt,bottom=3pt,sharp corners,title={#1},fonttitle=\bfseries}
\newtcolorbox{merke}{colback=gray!8,colframe=gray!60!black,boxrule=0.6pt,left=6pt,right=6pt,top=3pt,bottom=3pt,sharp corners}
\newcommand{\cb}{$\square$\ }
\newcommand{\teil}[1]{\par\medskip\textbf{#1)}\ }
